%% Question bank chapter file - generated by bank_tools.py build (do not edit while the bank is still being built:
%% it is regenerated from the per-test data in Question Bank\data). Every question is one bankq block:
%%   \begin{bankq}[id=..., source={...}, paper=.., qno=.., type=..., subtopic={...}, answer={...}, answer_src=key|solved,
%%                 working={...}, difficulty=1-10, diff_parts={K S I C F}, diff_note={...},
%%                 group=<passage id or empty>, also_in={duplicate copies in other tests}]
%% difficulty = min(10, K+S+I+C+F-1): K concepts 1-3, S steps 1-4, I insight 0-2, C computation 0-2, F format 0-1
%%   ...question LaTeX, options, figures (\bfig / \bfigi -> img/)...
%%   \end{bankq}
%% Passages shared by several questions are bankpassage blocks; their questions carry group=<passage id>.
\documentclass[10pt]{article}
\usepackage[a4paper,margin=1.8cm]{geometry}
\usepackage[T1]{fontenc}
\usepackage{amsmath,amssymb,mathtools}
\usepackage{graphicx}
\usepackage[export]{adjustbox}
\usepackage{xcolor,needspace,keyval}
\usepackage{enumitem}
\usepackage{booktabs,array,tabularx}
\usepackage[hidelinks]{hyperref}
\setlength{\parindent}{0pt}
\setlength{\parskip}{3pt}
\makeatletter
\define@key{bq}{id}{\def\bq@id{#1}}\define@key{bq}{source}{\def\bq@src{#1}}\define@key{bq}{paper}{\def\bq@paper{#1}}
\define@key{bq}{qno}{\def\bq@qno{#1}}\define@key{bq}{type}{\def\bq@type{#1}}\define@key{bq}{subtopic}{\def\bq@sub{#1}}
\define@key{bq}{answer}{\def\bq@ans{#1}}\define@key{bq}{answer_src}{\def\bq@asrc{#1}}\define@key{bq}{working}{\def\bq@work{#1}}
\define@key{bq}{group}{\def\bq@group{#1}}\define@key{bq}{also_in}{\def\bq@also{#1}}
\define@key{bq}{difficulty}{\def\bq@diff{#1}}\define@key{bq}{diff_parts}{\def\bq@dparts{#1}}\define@key{bq}{diff_note}{\def\bq@dnote{#1}}
\newenvironment{bankq}[1][]{\def\bq@work{}\def\bq@also{}\def\bq@group{}\def\bq@diff{?}\def\bq@dparts{}\def\bq@dnote{}\setkeys{bq}{#1}%
  \par\needspace{4\baselineskip}\medskip\noindent{\small\textbf{\bq@id}\hfill\color{gray}[\bq@src, Paper \bq@paper, Q\bq@qno]
  \bq@type\ $\cdot$ \bq@sub\ $\cdot$ \textbf{D\bq@diff}}\par\nopagebreak\relax}
  {\par{\small\color{teal}Answer: \textbf{\bq@ans}\ifx\bq@work\@empty\else\ (solved: \bq@work)\fi
  \ifx\bq@also\@empty\else\par\color{gray}Also in: \bq@also\fi}\par}
\define@key{bp}{id}{\def\bp@id{#1}}\define@key{bp}{source}{\def\bp@src{#1}}\define@key{bp}{questions}{\def\bp@qs{#1}}
\newenvironment{bankpassage}[1][]{\setkeys{bp}{#1}\par\needspace{4\baselineskip}\bigskip\noindent
  {\small\textbf{Passage \bp@id}\hfill\color{gray}[\bp@src; questions \bp@qs]}\par\nopagebreak\relax\itshape}{\par}
\makeatother
\newcommand{\bfig}[2]{\par\smallskip{\centering\includegraphics[width=#1pt,max width=0.62\linewidth,max height=0.32\textheight]{img/#2}\par}\smallskip}
\newcommand{\bfigi}[2]{\includegraphics[width=#1pt,max width=\linewidth,max height=0.32\textheight,valign=c]{img/#2}}
\newcommand{\bmissing}[1]{\textit{[figure missing in the source conversion: \detokenize{#1}]}}
\newenvironment{qoptions}{\begin{itemize}[leftmargin=1.8em,itemsep=1pt,topsep=2pt,label={}]}{\end{itemize}}
\newcommand{\opt}[2]{\item[(#1)] #2}
\providecommand{\fwstart}{}\providecommand{\fwend}{}\providecommand{\qfullsplit}{}
\newcolumntype{L}[1]{>{\raggedright\arraybackslash}p{#1}}
\newcolumntype{C}[1]{>{\centering\arraybackslash}p{#1}}
\newcolumntype{R}[1]{>{\raggedleft\arraybackslash}p{#1}}
\providecommand{\sectionhdr}[1]{}\providecommand{\subsectionhdr}[1]{}


\begin{document}
{\Large\bfseries Hydrocarbons --- Question Bank}\par\smallskip
{\small Chemistry chapter C17. Questions from all converted test papers, grouped by source series. Answers: printed key where it matches the paper, otherwise solved (answer\_src=solved, with working). Exact repeats across papers appear once, with also\_in listing the other copies.}\par\medskip\hrule

\section*{Score-I papers}
\begin{bankq}[id=SI12-P1-CHEM-S1-Q03, source={Score-I - Test 12 (09 Mar 2022)}, paper=1, qno=3, type=multi, subtopic={Alkenes}, answer={B,C}, answer_src=key, difficulty=2, diff_parts={K1 S1 I0 C0 F1}, diff_note={concepts: addition of halogens to alkenes | steps: options A and B; options C and D to identify the vicinal dichloride}]
Vicinal-dichloride is obtained as a product in reaction -
\begin{center}\small
\begin{tabular}{@{}l@{\ }l@{\qquad}l@{\ }c@{}}
(A) & CH$_3$--CH=O $\xrightarrow{\text{PCl}_5}$ & (B) & \bfigi{60.1}{SI12_p1_chem_s1_q03_B.png}\rule{0pt}{26pt}\\[8pt]
(C) & \bfigi{67.7}{SI12_p1_chem_s1_q03_C.png}\rule{0pt}{32pt} & (D) & \bfigi{68.1}{SI12_p1_chem_s1_q03_D.png}\rule{0pt}{40pt}
\end{tabular}
\end{center}
\end{bankq}
\begin{bankq}[id=SI12-P1-CHEM-S3-Q06, source={Score-I - Test 12 (09 Mar 2022)}, paper=1, qno=6, type=numerical, subtopic={Mechanisms}, answer={3}, answer_src=key, difficulty=5, diff_parts={K2 S3 I0 C0 F1}, diff_note={concepts: carbocation mechanism, degree of unsaturation | steps: first step of the sequence; second step; third step; products D and E; degree of unsaturation of each}]
Consider following sequence of reactions
\bfig{164.8}{SI12_p1_chem_s3_q06_1.png}
Total sum of degree of unsaturations of product D \& E will be
\end{bankq}
\begin{bankq}[id=SI12-P2-CHEM-S2-Q03, source={Score-I - Test 12 (09 Mar 2022)}, paper=2, qno=3, type=numerical, subtopic={Alkenes}, answer={2}, answer_src=solved, working={Only the symmetric alkenes cis- and trans-but-2-ene give 2-bromobutane either way}, difficulty=2, diff_parts={K1 S1 I0 C0 F1}, diff_note={concepts: regioselectivity of HBr addition with and without peroxide | steps: list the four C4H8 alkenes; only symmetric alkenes give the same product}]
Number of alkenes of formula C$_4$H$_8$ which give same product with HBr in presence as well as absence of peroxide are.
\end{bankq}

\section*{TOAS papers}
\begin{bankq}[id=TOAS02-P1-CHEM-S2-Q10, source={TOAS - Test 02 (13 Feb 2022)}, paper=1, qno=10, type=multi, subtopic={Alkenes}, answer={B,C}, answer_src=solved, working={A: 2,3-dimethylbutan-2-ol (rearranged, achiral); B: 3,3-dimethylbutan-1-ol (1$^\circ$, achiral); C: 3,3-dimethylbutan-2-ol (2$^\circ$, chiral); A and B differ in carbon skeleton, not positional isomers}, difficulty=4, diff_parts={K2 S2 I0 C0 F1}, diff_note={concepts: hydration of alkenes, carbocation rearrangement | steps: product A by acid hydration with rearrangement; product B by hydroboration; product C by oxymercuration; chirality and relations}]
\bfigi{153.7}{TOAS02_p1_chem_s2_q10_1.png}\par
Selelct correct statement form the following about major product A, B and C :
\begin{qoptions}
\opt{A}{A is $3^{\mathrm o}$ alcohol and show optical isomersim}
\opt{B}{B is $1^{\mathrm o}$ alcohol and does not show optical isomersim}
\opt{C}{C is $2^{\mathrm o}$ alcohol and shows optical isomersim}
\opt{D}{A and B are positional isomer of each other}
\end{qoptions}
\end{bankq}
\begin{bankq}[id=TOAS03-P1-CHEM-S1-Q02, source={TOAS - Test 03 (27 Feb 2022)}, paper=1, qno=2, type=single, subtopic={Alkenes}, answer={C}, answer_src=key, difficulty=1, diff_parts={K1 S1 I0 C0 F0}, diff_note={concepts: hydration of alkenes | steps: Markovnikov product in A and B; in C and D}]
Which of following gives $2^\circ$ alcohol major product?
\begin{qoptions}
\opt{A}{\bfigi{80.6}{TOAS03_p1_chem_s1_q02_A.png}}
\opt{B}{\bfigi{94.8}{TOAS03_p1_chem_s1_q02_B.png}}
\opt{C}{\bfigi{161.1}{TOAS03_p1_chem_s1_q02_C.png}}
\opt{D}{\bfigi{65.9}{TOAS03_p1_chem_s1_q02_D.png}}
\end{qoptions}
\end{bankq}
\begin{bankq}[id=TOAS03-P1-CHEM-S2-Q08, source={TOAS - Test 03 (27 Feb 2022)}, paper=1, qno=8, type=multi, subtopic={Mechanisms}, answer={B,C,D}, answer_src=key, difficulty=3, diff_parts={K1 S2 I0 C0 F1}, diff_note={concepts: kinetic versus thermodynamic control | steps: major product by the more stable carbocation; product at higher temperature; identify P and Q}]
\bfigi{125.6}{TOAS03_p1_chem_s2_q08_1.png}

Correct option(s) for above reaction sequence
\begin{qoptions}
\opt{A}{(P) is \bfigi{46.3}{TOAS03_p1_chem_s2_q08_2.png}}
\opt{B}{(Q) is \bfigi{46.3}{TOAS03_p1_chem_s2_q08_3.png}}
\opt{C}{On increasing the temperature \% of (Q) increases in product mixture}
\opt{D}{Product (P) is formed by more stable carbocation formation}
\end{qoptions}
\end{bankq}
\begin{bankq}[id=TOAS03-P1-CHEM-S3-Q02, source={TOAS - Test 03 (27 Feb 2022)}, paper=1, qno=2, type=numerical, subtopic={Alkenes}, answer={2}, answer_src=key, difficulty=2, diff_parts={K1 S1 I0 C0 F1}, diff_note={concepts: alkenes C5H10 and electrophilic addition | steps: pent-1-ene and pent-2-ene give 2-chloropentane; count constitutional possibilities}]
An unknown unbranched hydrocarbon with molecular formula $\mathrm{C_5H_{10}}$ decolourise $\mathrm{Br_2+H_2O}$ and on reaction with H--Cl/$\mathrm{CCl_4}$ give 2-chloro pentane as major product. Identify total number of possible reactants (excluding stereochemistry) -
\end{bankq}
\begin{bankq}[id=TOAS03-P2-CHEM-S1-Q01, source={TOAS - Test 03 (27 Feb 2022)}, paper=2, qno=1, type=multi, subtopic={Alkenes}, answer={A,B,D}, answer_src=key, difficulty=2, diff_parts={K1 S1 I0 C0 F1}, diff_note={concepts: Markovnikov's rule | steps: unsymmetrical alkenes in options A and B; in C and D}]
Markovnikoff's rule is applicable on which of the following alkene :
\begin{qoptions}
\opt{A}{\bfigi{43.0}{TOAS03_p2_chem_s1_q01_A.png}}
\opt{B}{\bfigi{38.4}{TOAS03_p2_chem_s1_q01_B.png}}
\opt{C}{\bfigi{38.2}{TOAS03_p2_chem_s1_q01_C.png}}
\opt{D}{\bfigi{38.7}{TOAS03_p2_chem_s1_q01_D.png}}
\end{qoptions}
\end{bankq}
\begin{bankq}[id=TOAS03-P2-CHEM-S1-Q02, source={TOAS - Test 03 (27 Feb 2022)}, paper=2, qno=2, type=multi, subtopic={Alkynes}, answer={A,C,D}, answer_src=key, difficulty=3, diff_parts={K1 S2 I0 C0 F1}, diff_note={concepts: reactions of alkynes | steps: hydration; addition of HX; reduction; check the products}]
Identify reactions correctly matched with their major products ?
\begin{qoptions}
\opt{A}{\bfigi{175.6}{TOAS03_p2_chem_s1_q02_A.png}}
\opt{B}{Me--C$\equiv$CH \bfigi{142.1}{TOAS03_p2_chem_s1_q02_2.png}}
\opt{C}{Me--C$\equiv$CH \bfigi{106.5}{TOAS03_p2_chem_s1_q02_3.png}}
\opt{D}{\bfigi{183.2}{TOAS03_p2_chem_s1_q02_D.png}}
\end{qoptions}
\end{bankq}
\begin{bankq}[id=TOAS03-P2-CHEM-S3-Q01, source={TOAS - Test 03 (27 Feb 2022)}, paper=2, qno=1, type=numerical, subtopic={Alkenes}, answer={1}, answer_src=key, difficulty=2, diff_parts={K1 S1 I0 C0 F1}, diff_note={concepts: stereospecific addition of bromine | steps: anti addition to trans-but-2-ene gives meso product only}]
How many product(s) is/are formed when trans bute-2-ene reacts with $\mathrm{Br_2}$ / $\mathrm{CCl_4}$
\end{bankq}
\begin{bankq}[id=TOAS04-P1-CHEM-S1-Q04, source={TOAS - Test 04 (15 Mar 2022)}, paper=1, qno=4, type=single, subtopic={Alkynes}, answer={D}, answer_src=key, difficulty=1, diff_parts={K1 S1 I0 C0 F0}, diff_note={concepts: alkylation of a terminal alkyne | steps: sodium acetylide formation; SN2 with ethyl iodide}]
What is the final product of the given reaction

$\mathrm{CH_3{-}C{\equiv}CH}\xrightarrow[\mathrm{(ii)\ CH_3{-}CH_2{-}I}]{\mathrm{(i)\ Na}}\mathrm{Product}$
\begin{qoptions}
\opt{A}{$\mathrm{CH_2{=}CH{-}CH_2{-}CH_3}$}
\opt{B}{$\mathrm{CH_3{-}CH_2{-}C{\equiv}C{-}CH_2{-}CH_3}$}
\opt{C}{$\mathrm{CH{\equiv}C{-}CH_2{-}CH_2{-}CH_3}$}
\opt{D}{$\mathrm{CH_3{-}C{\equiv}C{-}CH_2{-}CH_3}$}
\end{qoptions}
\end{bankq}
\begin{bankq}[id=TOAS04-P2-CHEM-S1-Q05, source={TOAS - Test 04 (15 Mar 2022)}, paper=2, qno=5, type=multi, subtopic={Alkynes}, answer={A,C,D}, answer_src=key, difficulty=2, diff_parts={K1 S1 I0 C0 F1}, diff_note={concepts: reactions of alkynes | steps: products for options A to D}]
Identify reactions correctly matched with their major products ?
\begin{qoptions}
\opt{A}{\bfigi{186.8}{TOAS04_p2_chem_s1_q05_A.png}}
\opt{B}{$\mathrm{Me{-}C{\equiv}CH}$ \bfigi{132.5}{TOAS04_p2_chem_s1_q05_2.png}}
\opt{C}{$\mathrm{Me{-}C{\equiv}CH}$ \bfigi{101.4}{TOAS04_p2_chem_s1_q05_3.png}}
\opt{D}{\bfigi{207.4}{TOAS04_p2_chem_s1_q05_D.png}}
\end{qoptions}
\end{bankq}
\begin{bankq}[id=TOAS05-P2-CHEM-S1-Q01, source={TOAS - Test 05 (03 Apr 2022)}, paper=2, qno=1, type=single, subtopic={Alkynes}, answer={B}, answer_src=key, difficulty=1, diff_parts={K1 S1 I0 C0 F0}, diff_note={concepts: functional group isomers from alkyne hydration | steps: hydration products of the options; keto-enol related isomers}]
\bfig{127.5}{TOAS05_p2_chem_s1_q01_1.png}
If $X$ and $Y$ are functional group isomers, then what is the structure of reactant (Alkyne).
\begin{qoptions}
\opt{A}{$\mathrm{CH_3{-}C{\equiv}C{-}CH_3}$}
\opt{B}{$\mathrm{CH_3{-}CH_2{-}C{\equiv}C{-}H}$}
\opt{C}{$\mathrm{CH_2{=}CH{-}CH{=}CH_2}$}
\opt{D}{$\mathrm{CH_3{-}CH{=}C{=}CH_2}$}
\end{qoptions}
\end{bankq}
\begin{bankq}[id=TOAS05-P2-CHEM-S1-Q05, source={TOAS - Test 05 (03 Apr 2022)}, paper=2, qno=5, type=single, subtopic={Alkenes}, answer={A}, answer_src=key, difficulty=1, diff_parts={K1 S1 I0 C0 F0}, diff_note={concepts: regioselectivity of HBr addition | steps: symmetric alkenes give the same product with and without peroxide}]
Which of the following alkene gives same product on reaction with $\mathrm{HBr/CCl_4}$ and $\mathrm{HBr/H_2O_2/h\nu}$
\begin{qoptions}
\opt{A}{\bfigi{40.7}{TOAS05_p2_chem_s1_q05_A.png}}
\opt{B}{\bfigi{30.4}{TOAS05_p2_chem_s1_q05_B.png}}
\opt{C}{\bfigi{43.0}{TOAS05_p2_chem_s1_q05_C.png}}
\opt{D}{\bfigi{41.9}{TOAS05_p2_chem_s1_q05_D.png}}
\end{qoptions}
\end{bankq}
\begin{bankq}[id=TOAS05-P2-CHEM-S2-Q02, source={TOAS - Test 05 (03 Apr 2022)}, paper=2, qno=2, type=numerical, subtopic={Alkenes}, answer={4}, answer_src=solved, working={X: 2-methylbut-2-ene (major) and 2-methylbut-1-ene; Br$_2$ gives 2,3-dibromo-2-methylbutane and 1,2-dibromo-2-methylbutane, each chiral (racemic): 4 products counting enantiomers}, difficulty=3, diff_parts={K1 S2 I0 C0 F1}, diff_note={concepts: bromination of an alkene mixture | steps: alkenes in X; dibromo product of each; chiral products as pairs}]
\bfigi{37.5}{TOAS05_p2_chem_s2_q02_1.png}\bfigi{113.1}{TOAS05_p2_chem_s2_q02_2.png}\par
Total number of possible dibromo products in final product mixture ($Y$)
\end{bankq}
\begin{bankq}[id=TOAS06-P1-CHEM-S2-Q09, source={TOAS - Test 06 (17 Apr 2022)}, paper=1, qno=9, type=multi, subtopic={Alkanes}, answer={A,C}, answer_src=key, difficulty=4, diff_parts={K1 S2 I1 C0 F1}, diff_note={concepts: monobromination of an alkane | steps: distinct hydrogens; chiral products; stereoisomers and enantiomeric pairs | I: count stereoisomers carefully}]
\bfigi{163.0}{TOAS06_p1_chem_s2_q08_2.png}
\begin{qoptions}
\opt{A}{Total 8 monobrominated product are form}
\opt{B}{Total 5 monobrominated product are form excluding stereoisomer}
\opt{C}{6 monobrominated are optically active}
\opt{D}{4 enantiomeric pair are form}
\end{qoptions}
\end{bankq}
\begin{bankq}[id=TOAS06-P1-CHEM-S3-Q03, source={TOAS - Test 06 (17 Apr 2022)}, paper=1, qno=3, type=numerical, subtopic={Alkenes}, answer={5}, answer_src=solved, working={NBS allylic bromination: Br at the tertiary ring C or the exocyclic CH2 (one radical) and at the ring CH2 or exocyclic CH2 (other radical): four constitutional products; the ring CH2 one has two stereocentres (cis/trans diastereomers separable), the others are single racemates or achiral: 5 fractions}, difficulty=4, diff_parts={K1 S2 I1 C0 F1}, diff_note={concepts: allylic bromination and fractions | steps: radical formation; products from each radical; stereoisomer pairs separable by distillation; count | I: count fractions not products}]
\bfigi{199.9}{TOAS06_p1_chem_s3_q03_1.png}\\
Find number of fraction obtained after fractional distillation of mono-brominated products
\end{bankq}
\begin{bankq}[id=TOAS07-P1-CHEM-S1-Q02, source={TOAS - Test 07 (01 May 2022)}, paper=1, qno=2, type=single, subtopic={Alkenes}, answer={D}, answer_src=key, difficulty=1, diff_parts={K1 S1 I0 C0 F0}, diff_note={concepts: addition reactions of alkenes | steps: Markovnikov and anti-Markovnikov addition; product in each option}]
Reaction correctly matched with its major product is :
\begin{qoptions}
\opt{A}{\bfigi{132.7}{TOAS07_p1_chem_s1_q02_A.png}}
\opt{B}{\bfigi{158.9}{TOAS07_p1_chem_s1_q02_B.png}}
\opt{C}{\bfigi{128.1}{TOAS07_p1_chem_s1_q02_C.png}}
\opt{D}{1-butene $\xrightarrow{\mathrm{HBr+H_2O_2}}$ 1-bromobutane}
\end{qoptions}
\end{bankq}
\begin{bankq}[id=TOAS07-P1-CHEM-S2-Q05, source={TOAS - Test 07 (01 May 2022)}, paper=1, qno=5, type=multi, subtopic={Alkenes}, answer={A,B,C,D}, answer_src=key, difficulty=3, diff_parts={K1 S2 I0 C0 F1}, diff_note={concepts: reactions of alkenes | steps: identify the reaction type; product in options A and B; options C and D}]
\bfigi{99.9}{TOAS07_p1_chem_s2_q05_1.png}
\begin{qoptions}
\opt{A}{\bfigi{38.7}{TOAS07_p1_chem_s2_q05_A.png}}
\opt{B}{\bfigi{46.0}{TOAS07_p1_chem_s2_q05_B.png}}
\opt{C}{\bfigi{39.2}{TOAS07_p1_chem_s2_q05_C.png}}
\opt{D}{\bfigi{46.1}{TOAS07_p1_chem_s2_q05_D.png}}
\end{qoptions}
\end{bankq}
\begin{bankpassage}[id=TOAS07-P1-CHEM-S3-G1-2, source={TOAS - Test 07 (01 May 2022)}, questions={1-2}]
Consider the reactions given below and answer the following questions
\bfig{86.5}{TOAS07_p1_chem_s3_para1_2_1.png}
(II) (R)--2--Chloropentane $\xrightarrow[h\nu]{\mathrm{Cl_2}}$ `Y'
\end{bankpassage}
\begin{bankq}[id=TOAS07-P1-CHEM-S3-Q01, source={TOAS - Test 07 (01 May 2022)}, paper=1, qno=1, type=numerical, subtopic={Alkanes}, answer={8.00}, answer_src=key, difficulty=3, diff_parts={K1 S2 I0 C0 F1}, diff_note={concepts: monochlorination of an alkane | steps: distinct hydrogen environments; stereo-centres created; count products}, group=TOAS07-P1-CHEM-S3-G1-2]
Find out total number of monochlorinated product `X'
\end{bankq}
\begin{bankq}[id=TOAS07-P1-CHEM-S3-Q02, source={TOAS - Test 07 (01 May 2022)}, paper=1, qno=2, type=numerical, subtopic={Alkanes}, answer={7.00}, answer_src=key, difficulty=5, diff_parts={K1 S2 I1 C1 F1}, diff_note={concepts: dichlorination and fractional distillation | steps: constitutional dichloro isomers; stereoisomers of each; diastereomers separate and enantiomers do not; count fractions | I: enantiomers stay together}, group=TOAS07-P1-CHEM-S3-G1-2]
If the mixture of dichlorinated products (Y) is subjected to fractional distillation, how many fractions would be obtained?
\end{bankq}
\begin{bankq}[id=TOAS08-P1-CHEM-S1-Q06, source={TOAS - Test 08 (15 May 2022)}, paper=1, qno=6, type=multi, subtopic={Alkenes}, answer={A,B,C,D}, answer_src=key, difficulty=2, diff_parts={K1 S1 I0 C0 F1}, diff_note={concepts: incorrect products of reactions | steps: reactions of options A to D; compare the claimed products}]
Which reaction(s) among the following is not represent correct product formation :
\begin{qoptions}
\opt{A}{\bfigi{114.8}{TOAS08_p1_chem_s1_q06_A.png}}
\opt{B}{\bfigi{115.3}{TOAS08_p1_chem_s1_q06_B.png}}
\opt{C}{\bfigi{160.6}{TOAS08_p1_chem_s1_q06_C.png}}
\opt{D}{\bfigi{157.3}{TOAS08_p1_chem_s1_q06_D.png}}
\end{qoptions}
\end{bankq}
\begin{bankq}[id=TOAS08-P2-CHEM-S1-Q06, source={TOAS - Test 08 (15 May 2022)}, paper=2, qno=6, type=multi, subtopic={Alkenes}, answer={B,C}, answer_src=key, difficulty=4, diff_parts={K1 S2 I1 C0 F1}, diff_note={concepts: NBS reactions and stereoisomers | steps: products P, Q, R; allylic substitution counts; stereoisomerism | I: allylic radical intermediates}]
Choose the correct statement(s) regarding following reaction :
\bfig{184.7}{TOAS08_p2_chem_s1_q06_1.png}
\begin{qoptions}
\opt{A}{$(Q)$, will give only 1 product with NBS}
\opt{B}{$(R)$, is optically inactive due to internal compensation}
\opt{C}{Product $(Q)$ is trans-but-2-ene}
\opt{D}{$(P)$, can exist in four stereoisomeric forms.}
\end{qoptions}
\end{bankq}
\begin{bankq}[id=TOAS09-P1-CHEM-S1-Q01, source={TOAS - Test 09 (30 May 2022)}, paper=1, qno=1, type=single, subtopic={Alkenes}, answer={C}, answer_src=key, difficulty=1, diff_parts={K1 S1 I0 C0 F0}, diff_note={concepts: stereochemistry of alkene addition | steps: syn or anti addition for each alkene; meso product}]
Which of the following will give non resolvable (meso) product :
\begin{qoptions}
\opt{A}{\bfigi{75.4}{TOAS09_p1_chem_s1_q01_A.png}}
\opt{B}{\bfigi{91.1}{TOAS09_p1_chem_s1_q01_B.png}}
\opt{C}{\bfigi{84.0}{TOAS09_p1_chem_s1_q01_C.png}}
\opt{D}{\bfigi{90.1}{TOAS09_p1_chem_s1_q01_D.png}}
\end{qoptions}
\end{bankq}
\begin{bankq}[id=TOAS09-P1-CHEM-S1-Q02, source={TOAS - Test 09 (30 May 2022)}, paper=1, qno=2, type=single, subtopic={Alkynes}, answer={D}, answer_src=key, difficulty=2, diff_parts={K1 S2 I0 C0 F0}, diff_note={concepts: hydration and reduction sequence | steps: alkyne hydration to ketone; reduction to alcohol; compare the options}]
Which compound is suitable for `P' in following reaction :
\par\smallskip P $\xrightarrow[\mathrm{H_2SO_4}]{\mathrm{HgSO_4}}\ \xrightarrow[\mathrm{H_2O}]{\mathrm{NaBH_4}}$ \bfigi{108.7}{TOAS09_p1_chem_s1_q02_1.png}
\begin{qoptions}
\opt{A}{\bfigi{62.5}{TOAS09_p1_chem_s1_q02_A.png}}
\opt{B}{\bfigi{62.5}{TOAS09_p1_chem_s1_q02_B.png}}
\opt{C}{\bfigi{60.4}{TOAS09_p1_chem_s1_q02_C.png}}
\opt{D}{\bfigi{47.3}{TOAS09_p1_chem_s1_q02_D.png}}
\end{qoptions}
\end{bankq}
\begin{bankq}[id=TOAS09-P1-CHEM-S2-Q08, source={TOAS - Test 09 (30 May 2022)}, paper=1, qno=8, type=multi, subtopic={Alkenes}, answer={A,B}, answer_src=key, difficulty=2, diff_parts={K1 S1 I0 C0 F1}, diff_note={concepts: oxidative cleavage of alkenes | steps: products with KMnO4 and with ozonolysis for each alkene; compare}]
Which of the following will give same product with H$^{\oplus}$, KMnO$_4$, $\Delta$ and O$_3$ / Zn, H$_2$O ?
\begin{qoptions}
\opt{A}{\bfigi{35.8}{TOAS09_p1_chem_s2_q08_A.png}}
\opt{B}{\bfigi{48.7}{TOAS09_p1_chem_s2_q08_B.png}}
\opt{C}{\bfigi{37.7}{TOAS09_p1_chem_s2_q08_C.png}}
\opt{D}{\bfigi{28.2}{TOAS09_p1_chem_s2_q08_D.png}}
\end{qoptions}
\end{bankq}
\begin{bankpassage}[id=TOAS09-P1-CHEM-S3-G5-6, source={TOAS - Test 09 (30 May 2022)}, questions={5-6}]
Among the following alkenes :
\bfig{318.2}{TOAS09_p1_chem_s3_para5_6_1.png}
\end{bankpassage}
\begin{bankq}[id=TOAS09-P1-CHEM-S3-Q05, source={TOAS - Test 09 (30 May 2022)}, paper=1, qno=5, type=numerical, subtopic={Alkenes}, answer={3.00}, answer_src=key, difficulty=4, diff_parts={K1 S2 I1 C0 F1}, diff_note={concepts: ozonolysis of alkenes | steps: products of ozonolysis for each alkene; which are meso; count | I: cyclic alkenes give meso dicarbonyls}, group=TOAS09-P1-CHEM-S3-G5-6]
How many alkenes which give meso compound on ozonolysis in presence of Zn?
\end{bankq}
\begin{bankq}[id=TOAS09-P1-CHEM-S3-Q06, source={TOAS - Test 09 (30 May 2022)}, paper=1, qno=6, type=numerical, subtopic={Alkenes}, answer={5.00}, answer_src=key, difficulty=4, diff_parts={K1 S2 I1 C0 F1}, diff_note={concepts: ozonolysis and optical activity | steps: products for each alkene; chirality of each product; count | I: chirality of the carbonyl products}, group=TOAS09-P1-CHEM-S3-G5-6]
From the following options identify alkenes which give optically active compound as one of the product on ozonolysis in presence of Zn.
\end{bankq}
\begin{bankq}[id=TOAS09-P2-CHEM-S1-Q09, source={TOAS - Test 09 (30 May 2022)}, paper=2, qno=9, type=single, subtopic={Alkynes}, answer={C}, answer_src=key, difficulty=1, diff_parts={K1 S1 I0 C0 F0}, diff_note={concepts: reduction of alkynes | steps: reagent and product; Na/liquid ammonia versus catalytic hydrogenation}]
\bfig{145.8}{TOAS09_p2_chem_s1_q09_1.png}
Major product of above reaction :
\begin{qoptions}
\opt{A}{\bfigi{21.6}{TOAS09_p2_chem_s1_q09_A.png}}
\opt{B}{CH$_3$--CH$_3$}
\opt{C}{CH$_2$=CH$_2$}
\opt{D}{HC$\equiv$CH}
\end{qoptions}
\end{bankq}
\end{document}
